\section{Phương pháp nghiên cứu}
\label{method}

\subsection{Dữ liệu}

Nghiên cứu sử dụng tập dữ liệu \textbf{FoodSeg103}, bao gồm \textbf{7.118 hình ảnh thực phẩm} được gán nhãn phân đoạn ở mức pixel cho \textbf{103 loại nguyên liệu khác nhau}. Dữ liệu được trích xuất từ tập Recipe1M và được gán nhãn thủ công. Tập dữ liệu được chia thành:
\begin{itemize}
    \item 4.983 ảnh dùng để huấn luyện;
    \item 2.135 ảnh dùng để kiểm tra hoặc đánh giá mô hình.
\end{itemize}

Các bước tiền xử lý dữ liệu bao gồm: resize ảnh, random crop, horizontal flip, color jitter. Những kỹ thuật này giúp cải thiện khả năng tổng quát hóa của mô hình.

\subsection{Mô hình cơ sở}

Mô hình baseline sử dụng trong nghiên cứu là \textbf{DeepLabV3}. DeepLabV3 sử dụng atrous convolution để điều chỉnh trường nhìn (field-of-view) của mạng và mô-đun ASPP để tổng hợp thông tin ngữ cảnh đa tỉ lệ. Các cấu hình sau sẽ được thử nghiệm:
\begin{itemize}
    \item DeepLabV3 + ResNet (baseline);
    \item DeepLabV3 + MobileNetV2;
    \item DeepLabV3 + MobileNetV3.
\end{itemize}
Mục tiêu là so sánh hiệu năng của các backbone này.

\subsection{Các cải tiến đề xuất}

\subsubsection{Boundary refinement}

Do các nguyên liệu trong ảnh thực phẩm thường có ranh giới không rõ ràng, nghiên cứu sẽ thử nghiệm các phương pháp cải thiện ranh giới phân đoạn như: sử dụng decoder kiểu DeepLabV3+; thêm các phương pháp tối ưu ranh giới đối tượng. Những phương pháp này giúp cải thiện độ chính xác ở vùng biên của các đối tượng trong ảnh.

\subsubsection{Test-time augmentation}

Một số kỹ thuật tăng cường dữ liệu tại thời điểm suy luận cũng sẽ được thử nghiệm, ví dụ: lật ảnh (flip), thay đổi tỷ lệ (scale). Các dự đoán từ nhiều biến đổi sẽ được kết hợp để cải thiện độ chính xác của mô hình.

\subsubsection{Tối ưu cho thiết bị di động}

Mô hình sau khi huấn luyện sẽ được tối ưu hóa để triển khai trên thiết bị di động. Các bước có thể bao gồm: xuất mô hình sang định dạng triển khai; tối ưu hóa suy luận; thử nghiệm quantization để giảm kích thước mô hình.

\subsection{Phương pháp đánh giá}

Hiệu năng của các mô hình sẽ được đánh giá bằng các chỉ số sau:
\begin{itemize}
    \item \textbf{Mean Intersection over Union (mIoU)};
    \item \textbf{Pixel Accuracy};
    \item \textbf{Inference time (latency)};
    \item \textbf{Kích thước mô hình}.
\end{itemize}
Các chỉ số này giúp đánh giá sự cân bằng giữa độ chính xác và hiệu quả tính toán của mô hình.
